\chapter{Notation and glossary}\label{ch:notation}\label{I:sec:glossary}\label{II:sec:glossary}\label{III:sec:glossary}
This chapter collects the notation and conventions of the book. Sections~\ref{gl:labels}--\ref{gl:named} hold for every chapter; they include the status labels (\S\ref{gl:labels}) and the cocycle normalization (\S\ref{gl:cocycle}). The remaining sections list the notation specific to individual chapters, with pointers to the definitions. The chapter proves nothing and can be consulted at any point.

\section{Status labels}\label{gl:labels}
Every statement in the series carries one of the following labels.
A label is never upgraded without a proof, and every change of label
is recorded in the revision log of the draft in which it is made.
\begin{description}
\item[\textbf{[P]}] \emph{Proved here}: a complete argument is
included in the Part in which the label appears, possibly using
results labeled \textbf{[E]}.  A statement proved under an
explicitly stated hypothesis is labeled \textbf{[P]} together with
the hypothesis, as in ``\textbf{[P]} under $(\mathrm H_D)$'', and the
hypothesis itself is listed separately with its own label.
\item[\textbf{[E]}] \emph{Established elsewhere}: a precisely cited
theorem from the literature or from an earlier Part, used as a black
box.  ``\textbf{[P]} (using \textbf{[E]})'' marks a proof that rests
on such inputs; ``\textbf{[E]}/\textbf{[P]}'' marks a statement that
is partly cited and partly proved.
\item[\textbf{[O]}] \emph{Proof outlined}: the steps are identified,
but at least one of them remains to be written out or checked.
\item[\textbf{[C]}] \emph{Conjectural}: new mathematical input is
required.  This includes the goals deferred to later Parts and the
hypotheses under which conditional statements are proved.
\item[\textbf{[P$^{\mathrm c}$]}] \emph{Proved modulo an exhaustive
exact computation}: the proof consists of hand arguments together
with a computation that is exhaustive over an explicitly described
finite set and uses exact arithmetic only (no sampling, and no
floating point without exact re-verification).  Each such script is
committed to the repository, cited by commit hash, and described in
an appendix and in the scripts' README.  At present this label is
used only in Chapter~\ref{ch:W}.
\end{description}

\section{Global data}\label{gl:global}
\begin{description}
\item[$K$, $n$, $\mu_n$] $K$ is a totally imaginary number field
containing the group $\mu_n$ of $n$-th roots of unity, and $n\ge1$ is
odd.  All archimedean places of $K$ are complex.
\item[$r$, $D$, $L$, $S$] $r$ is a prime with $(n,r)=1$; $D$ is a
central division algebra over $K$ of dimension $r^2$; $L\subset D$ is
a cyclic maximal subfield, so $[L:K]=r$; $S$ is the (finite) set of
places at which $D$ ramifies.  Since $r$ is prime, every $v\in S$ is
finite and $D_v$ is a division algebra; $|S|\ge2$, and $|S|$ is even
if $r=2$ (\S\ref{I:subsec:setting}).
\item[$r'$, $c$, $c_v$] Integers with $rr'\equiv1\pmod n$ and
$n\mid r-1+2c$; then $c_v:=r'c$ satisfies $n\mid r-1+2rc_v$ (\cref{I:ass:standing}).  Despite the subscript, $c_v$ does not depend on the
place $v$.  The parameter $c$ occurs \emph{before} the relabeling of
$\mu_n$ described in \S\ref{gl:cocycle}, and $c_v$ \emph{after} it.
\item[$\mathbb A_K$, $L_v$] $\mathbb A_K$ is the adele ring of $K$
(written $\mathbb A$ when no confusion is possible), and
$L_v:=L\otimes_KK_v=\prod_{w\mid v}L_w$.
\item[$G_D$] The algebraic group $D^\times$ over $K$, so that
$G_D(K_v)=D_v^\times$.
\item[Running example] $K=\mathbb Q(\zeta_3)$, $n=3$, $r=2$,
$D=(-1,4-\zeta_3)_K$, ramified exactly at
$S=\{\mathfrak p_3,\mathfrak p_7\}$ with $\mathfrak p_3=(1-\zeta_3)$
and $\mathfrak p_7=(3+\zeta_3)$; $L=K(i)=\mathbb Q(\zeta_{12})$,
$c=1$, $r'=2$, $c_v=2$.  The place $\mathfrak p_3$ divides $n$ and
$\mathfrak p_7$ does not.  The note
\texttt{notes/mjl\_running\_example.tex} works the example out across
the series.
\end{description}

\section{Local fields and absolute values}\label{gl:local}
\begin{description}
\item[$F$, $\mathcal O_F$, $\varpi$, $q_F$, $p$] $F$ is a
nonarchimedean local field of characteristic zero containing $\mu_n$
(usually $F=K_v$), with ring of integers $\mathcal O_F$, uniformizer
$\varpi$, residue field of order $q_F$, and residual characteristic
$p$.
\item[$|\cdot|_F$, $|\cdot|_p$] $|\cdot|_F$ is the normalized absolute
value, $|\varpi|_F=q_F^{-1}$; on $\mathbb Q_p$ it is written
$|\cdot|_p$, with $|p|_p=p^{-1}$.  In particular $|n|_F=1$ if
$p\nmid n$, and $|F^\times/F^{\times n}|=n^2|n|_F^{-1}$ (Chapter~\ref{ch:I},
proof of Proposition~5.5).
\item[Places dividing $n$; tame and wild places] A place $v$ (or the
field $F=K_v$) \emph{divides $n$} if $p\mid n$.  The places $v\in S$
with $p\nmid n$ are called \emph{tame}, and those with $p\mid n$
\emph{wild}.  In the running example $\mathfrak p_7$ is tame and
$\mathfrak p_3$ is wild ($|3|_{\mathfrak p_3}=3^{-2}$).
\end{description}

\section{Hilbert symbols}\label{gl:hilbert}
\begin{description}
\item[$(a,b)_F$] The $n$-th Hilbert symbol of $F$, with values in
$\mu_n$; Chapter~\ref{ch:IV} and the companion paper also write $(a,b)_{n,F}$ or
$(a,b)_n$.  It is bimultiplicative, nondegenerate and antisymmetric,
$(a,b)_F=(b,a)_F^{-1}$.  At complex places it is trivial.
\item[Norm compatibility] For a finite extension $F'/F$, $a\in
F^\times$ and $b\in F'^\times$, $(a,b)_{F'}=(a,N_{F'/F}b)_F$; in
particular $(a,b)_{F'}=(a,b)_F^{[F':F]}$ for $a,b\in F^\times$
[Serre, \emph{Local Fields}, Ch.~XIV].
\item[$(a,b)_{\mathbb A}$] For a number field $M\supseteq K$,
$(a,b)_{\mathbb A_M}=\prod_w(a_w,b_w)_{M_w}$.  Hilbert reciprocity:
$(a,b)_{\mathbb A_M}=1$ for $a,b\in M^\times$.
\end{description}

\section{Covers and the cocycle normalization}\label{gl:cocycle}
\begin{description}
\item[Cover, $p$, $\overline{H'}$] An \emph{$n$-fold cover} of a
locally compact group $H$ is a topological central extension
$1\to\mu_n\to\overline H\xrightarrow{\ p\ }H\to1$.  For a subgroup
$H'\subset H$ we write $\overline{H'}:=p^{-1}(H')$.  Chapters~\ref{ch:I} and~\ref{ch:II}
write $\overline H$; Parts~III, IV and the companion paper write
$\widetilde H$, $\widetilde Z$.
\item[Commutators] For lifts $\tilde x,\tilde y$ of commuting
elements, $[x,y]:=\tilde x\tilde y\tilde x^{-1}\tilde y^{-1}\in\mu_n$;
it does not depend on the lifts, and it is bimultiplicative on a
torus [Flicker--Kazhdan, \S3].
\item[Kazhdan--Patterson cocycle $\sigma_F$] The cocycle of
[Kazhdan--Patterson, \S0.1] and [Flicker--Kazhdan, \S2] on
$\mathrm{GL}(r,F)$ with twisting parameter $0$.  On the diagonal
torus,
$\sigma_F(\mathrm{diag}(a_i),\mathrm{diag}(b_i))=\prod_{i<j}(a_i,b_j)_F$.
\item[Parameter-$c'$ cover $\overline G^{(c')}$] The cover of
$\mathrm{GL}(r,F)$ with cocycle
$\sigma_F(g,h)\,(\det g,\det h)_F^{\,c'}$.  The comparison uses
$c'=c_v$; the parameter is forced (\cref{I:rem:forced}).  On the
diagonal torus the commutator of $\mathrm{diag}(t_i)$ and
$\mathrm{diag}(u_j)$ is $\prod_{i,j}(t_i,u_j)_F^{A_{ij}}$ with
$A=-I_r+(1+2c_v)J_r$, where $J_r$ is the all-ones matrix (\S\ref{IV:subsec:restriction}).
\item[Ambient cover $\overline{\mathcal G}_v$] The cover of
$\mathcal G_v=\mathrm{GL}(r,L_v)=\prod_{w\mid v}\mathrm{GL}(r,L_w)$
with cocycle
$\prod_{w\mid v}\sigma_{L_w}(g_w,h_w)\,(\det g_w,\det h_w)_{L_w}^{\,c_v}$,
the copies of $\mu_n$ being identified (\S\ref{I:subsec:ambient} and
equation~(1)).  Both covers compared in the series are subgroups of
it, so their common central subgroup $p^{-1}(K_v^\times)$ is literally
the same subgroup on the two sides.
\item[Restriction and relabeling] Pulled back along
$\iota_D:D_v^\times\hookrightarrow\mathcal G_v$, the ambient cocycle
is the local component of the cocycle
$\beta=\sigma_L\cdot(\mathrm{Nrd},\mathrm{Nrd})^c$ of the corrected
1988 paper [JC, equation~(4)].  Pulled back along the diagonal
embedding of $\mathrm{GL}(r,K_v)$, it is isomorphic, \emph{as an
extension}, to the extension with cocycle
$\sigma_F^{\,r}(\det,\det)_F^{\,c}$ (\cref{I:prop:restriction}).  The
group $\mu_n$ is relabeled \emph{once}, at every place and globally,
by $\zeta\mapsto\zeta^{r'}$; after relabeling the second extension is
$\overline G^{(c_v)}$.  Formulas state whether they hold before or
after relabeling.  The isomorphism is one of extensions, not an
equality of cocycles, and it is unique up to
$\tilde g\mapsto\chi(\det p(\tilde g))\,\tilde g$ with
$\chi:F^\times\to\mu_n$ continuous.  Central characters are therefore
never transported through it (\cref{I:rem:chidet}).
\item[Scalar centrality] Because $n\mid r-1+2rc_v$, the whole inverse
image $\overline Z=p^{-1}(F^\times)$ of the scalars is central
(\cref{I:lem:center}).  A cover with this property is called
\emph{scalar-central}.
\item[Global normalization] Globally one uses the normalized cocycle
$\tau_w(g,h)=\sigma_w(g,h)\,\kappa_w(g)\kappa_w(h)\kappa_w(gh)^{-1}$
of [Flicker--Kazhdan, \S2], where $\kappa_w$ is trivial for $w\mid
n\infty$.  The compact splitting is $\kappa^*_w(k)=(k,1)_\tau$ on
$\mathrm{GL}(r,\mathcal O_{L_w})$ for $w\nmid n\infty$, and the
rational splitting is $\iota(x)=s_\tau(x)\,\kappa_L(x)^{-1}$; in the
unnormalized models the rational splitting is $x\mapsto(x,1)_\sigma$
at every place (\cref{I:lem:splittings}).  The global ambient cover
$\overline{\mathcal G}(\mathbb A_L)$ has cocycle
$\tau_L\cdot(\det,\det)_{\mathbb A_L}^{\,c_v}$.
\item[Star map] For $h$ in the base group and any lift $\tilde h$,
$h^*:=\tilde h^{\,n}$.  It does not depend on the lift, it is
conjugation equivariant, $(ah)^*=a^*h^*$ for scalars $a$, and it is
multiplicative on every torus because $n$ is odd (\cref{III:def:star}; \S\ref{IV:subsec:star}).  For $\zeta\in\mu_n(F)$ one has
$\zeta^*=1$.
\end{description}

\section{Genuine objects, central characters and good elements}\label{gl:genuine}
\begin{description}
\item[Genuine representation] A representation of $\overline H$ on
which $\mu_n$ acts by the tautological character $\zeta\mapsto\zeta$.
\item[$\tilde\theta$, $\theta$] A genuine character $\tilde\theta$ of
$\overline Z$ (unitary from \S\ref{I:subsec:globalcovers} on; globally, a character of
$\overline Z(\mathbb A_K)$ trivial on $\iota(K^\times)$).  The same
$\tilde\theta$ serves both covers, since $\overline Z$ is a common
subgroup.  The ordinary character attached to it is
$\theta(a):=\tilde\theta(a^*)$; it is trivial on $\mu_n$, and
$\theta$ does not determine $\tilde\theta$ (\S\ref{IV:subsec:star}).
Chapter~\ref{ch:III} requires only that $\tilde\theta|_{\mu_n}$ be faithful.
\item[Genuine test functions] $C_c^\infty(\overline H,\tilde\theta)_{\mathrm{gen}}$
is the space of smooth functions with support compact modulo
$\overline Z$ and $f(z\tilde h)=\tilde\theta(z)^{-1}f(\tilde h)$ for
$z\in\overline Z$; in particular $f(\zeta\tilde h)=\zeta^{-1}f(\tilde h)$.
Characters and invariant class functions of genuine representations
transform by $\tilde\theta$ itself (\S\ref{III:sec:setting}).
\item[Regular, elliptic] An element of $\mathrm{GL}(r,F)$ is
\emph{regular} if its eigenvalues are distinct, and \emph{elliptic
regular} if moreover its characteristic polynomial is irreducible.
Since $r$ is prime, the regular elements of $D_v^\times$ are the
noncentral ones, and they are all elliptic.  An element of a cover is
regular (elliptic) if its image is.
\item[Good element] A regular $\tilde\gamma$ is \emph{good} if it
commutes with $p^{-1}(T_\gamma)$, where $T_\gamma$ is the centralizer
of $\gamma=p(\tilde\gamma)$; equivalently, $\tilde\gamma$ is central
in the covered torus $\overline{T_\gamma}$.  Genuine invariant
functions vanish at regular elements that are not good.  For a
maximal subfield $E$ of $D_v$ (or an elliptic torus $E^\times$ of
$\mathrm{GL}(r,F)$), the good elements of $\overline{E^\times}$ are
$p^{-1}(F^\times E^{\times n})=\overline Z\cdot\{t^*\}$, of index
$|n|_F^{-(r-1)}$, and every genuine irreducible representation of
$\overline{E^\times}$ has dimension $|n|_F^{-(r-1)/2}$ (\cref{I:prop:good}).  Covered elliptic tori are abelian exactly at the
places not dividing $n$.
\end{description}

\section{Orbital integrals and matching}\label{gl:matching}
\begin{description}
\item[$\Delta$] For regular semisimple $\gamma$,
$\Delta(\gamma)=|D(\gamma)|_F^{1/2}|\det\gamma|_F^{-(r-1)/2}$, where
$D(\gamma)$ is the discriminant of the (reduced) characteristic
polynomial and $\det=\mathrm{Nrd}$ on the division-algebra side.
Thus $\Delta(\gamma)^2$ is the Weyl discriminant
$|\det(1-\mathrm{Ad}(\gamma^{-1}))|_{\mathfrak g/\mathfrak t}|_F$
(\cref{I:def:oi} and \cref{I:lem:traceOI}).
\item[$\Phi$, $F=\Delta\Phi$] $\Phi(\tilde\gamma,\tilde f)$ is the
orbital integral over $T_\gamma\backslash G$ of a genuine test
function at a good regular $\tilde\gamma$, and
$F(\tilde\gamma,\tilde f)=\Delta(\gamma)\Phi(\tilde\gamma,\tilde f)$
is the normalized orbital integral, set to $0$ at regular elements
that are not good (\cref{I:def:oi}).  Chapter~\ref{ch:III} writes
$J(t,f)=\Delta(t)\,O(t,f)$ for the same function, with its own
measure conventions.
\item[Canonical identification $\jmath$] At $v\in S$ and for a maximal
subfield $E\subset D_v$ with an embedding $j:E\hookrightarrow
M_r(F)$, every isomorphism of covered tori that is the identity on
$\overline Z$ has the same restriction $\jmath_{E,j}$ to the good
elements; these restrictions induce a bijection $\jmath$ from the good
regular classes of the cover of $D_v^\times$ to the good regular
elliptic classes of the cover of $\mathrm{GL}(r,F)$ (\cref{I:thm:ident} and \cref{I:cor:classes}).
\item[Matching] Genuine test functions $\tilde f_D$ and $\tilde f$
\emph{match} if (M1)
$F(\tilde\gamma,\tilde f_D)=(-1)^{r-1}F(\jmath(\tilde\gamma),\tilde f)$
at every good regular $\tilde\gamma$, and (M2) $F(\tilde\delta,\tilde
f)=0$ at every good regular non-elliptic $\tilde\delta$ (\cref{I:def:matching}).  The sign sits in the matching of functions, so that
the expected character relation is
$\Theta_{\overline\pi_D}=(-1)^{r-1}\Theta_{\overline\pi}\circ\jmath$.
\item[Supercusp form] A genuine test function on the cover of
$\mathrm{GL}(r,F)$ whose integrals along the unipotent radicals of
all proper parabolic subgroups vanish, for example a matrix
coefficient of a supercuspidal representation [Flicker--Kazhdan,
\S13].
\item[Germ] The asymptotic behavior of orbital integrals near a
singular element, expressed through Shalika germs [Flicker--Kazhdan,
\S8]; used near the center in Chapter~\ref{ch:II}.
\end{description}

\section{Haar measures}\label{gl:measures}
\begin{description}
\item[Chapter~\ref{ch:I}] At $v\in S$ the Haar measures on $D_v^\times/F^\times$
and $\mathrm{GL}(r,F)/F^\times$ are arbitrary but fixed, and matching
is defined with respect to them; torus measures are transported by
conjugation and do not affect matching.  At $v\notin S$ the measures
correspond under the identifications $\psi_v$ (\S\S\ref{I:subsec:oi}, \ref{I:subsec:measures}).
\item[Chapter~\ref{ch:II}] At $v\in S$,
$\mathrm{vol}(D_v^\times/F^\times)=d(\mathrm{St})^{-1}$, where
$d(\mathrm{St})$ is the formal degree of the Steinberg representation
for the chosen measure on $\mathrm{GL}(r,F)/F^\times$ (\cref{II:def:measures}).
\item[Chapter~\ref{ch:III}] For the transform $L_D$, the quotients $H/Z$ and
$T/Z$ carry their Haar \emph{probability} measures (\S\ref{III:sec:local}).
\item[Formal degree] For a representation $\sigma$ that is square
integrable modulo the center, the constant $d(\sigma)$ in the Schur
orthogonality relations; it depends on the Haar measure.
\end{description}

\section{The division-side transform}\label{gl:transform}
\begin{description}
\item[$H$, $\widetilde H$, $L_D$] For a local division algebra $D$ of
degree $r$ over $F$, $H=D^\times$ and $\widetilde H$ is its cover from
Chapter~\ref{ch:I}.  The transform $L_D$ sends $\theta$-equivariant invariant
$L^2$ functions $\varphi$ on $H$ to genuine
$\tilde\theta$-equivariant invariant functions on $\widetilde H$ by
$(L_D\varphi)(z\,t^*)=b_D\,\tilde\theta(z)\,\Delta(t)\Delta(t^n)^{-1}\varphi(t)$
at good regular elements, and $0$ at the others; it is unitary
(\cref{III:def:L} and \cref{III:thm:unitary}).
\item[$b_D$, $c_D$] $b_D=|n|_F^{-(r-1)/2}$ and
$c_D=|n|_F^{-(r^2-1)/2}$ (\S\S\ref{III:sec:local}--\ref{III:sec:characters}).  Both equal $1$ at a
tame place; in the running example $b_D=3$ and $c_D=27$ at
$\mathfrak p_3$.
\item[Integrality of $L_D$] The question whether $L_D$ carries
irreducible characters to irreducible genuine characters.  It holds at
tame places (Chapter~\ref{ch:III}) and is open in general at wild places
(Chapter~\ref{ch:IV}; companion paper).
\end{description}

\section{Named results and hypotheses}\label{gl:named}
\begin{description}
\item[Theorems A, B, C] Chapter~\ref{ch:I}: the comparison of covered elliptic
tori, the matching of regular-supported test functions, and the
global transfer under $(\mathrm H_D)$ (with its converse under (H)).
\item[Theorems D, E, C$'$] Chapter~\ref{ch:II}: the equivariant section near the
center and its compatibility with $\jmath$; the matching of arbitrary
genuine test functions; and the global transfer under the weaker
hypothesis $(\mathrm H_D')$ (with its converse under $(\mathrm H')$).
\item[$(\mathrm H_D)$, (H), $(\mathrm H_D')$, $(\mathrm H')$]
Non-cancellation hypotheses: the character of a near-equivalence
packet at the places of $S$ is not identically zero on the relevant
matching test functions (\cref{I:def:HD} and
\cref{I:thm:converse}; \cref{II:def:HDprime} and \S\ref{II:subsec:Cprime}).  $(\mathrm H_D)$
holds when the $D$-side packet is finite (\cref{I:prop:HDfinite}); sufficient conditions for (H) are in \cref{I:rem:H}; and $(\mathrm H_D)\Rightarrow(\mathrm H_D')$,
$(\mathrm H)\Rightarrow(\mathrm H')$ (\S\ref{II:subsec:Cprime}).
\item[(D1)--(D4)] The steps left open by Chapter~\ref{ch:I}: matching of
arbitrary functions, local character relations, finiteness of
near-equivalence packets, and composition with [JC] (\S\ref{I:sec:deferred}).
Their status in the later Parts is recorded in \S\ref{I:subsec:partII}.
\item[Targets A, B, C] The companion paper's three local problems at
the wild place of the running example and beyond (Chapter~\ref{ch:W}, \S\ref{W:sec:intro}).
\end{description}

\begingroup
% macros of Part~I, Draft~11
\providecommand{\ov}{}\renewcommand{\ov}[1]{\overline{#1}}
\providecommand{\eps}{}\renewcommand{\eps}{\epsilon}
\providecommand{\A}{}\renewcommand{\A}{\mathbb{A}}
\providecommand{\C}{}\renewcommand{\C}{\mathbb{C}}
\providecommand{\Q}{}\renewcommand{\Q}{\mathbb{Q}}
\providecommand{\Z}{}\renewcommand{\Z}{\mathbb{Z}}
\providecommand{\cO}{}\renewcommand{\cO}{\mathcal{O}}
\providecommand{\gen}{}\renewcommand{\gen}{\mathrm{gen}}
\providecommand{\reg}{}\renewcommand{\reg}{\mathrm{reg}}
\providecommand{\tr}{}\renewcommand{\tr}{\operatorname{tr}}
\providecommand{\Gal}{}\renewcommand{\Gal}{\operatorname{Gal}}
\providecommand{\Aut}{}\renewcommand{\Aut}{\operatorname{Aut}}
\providecommand{\Nrd}{}\renewcommand{\Nrd}{\operatorname{Nrd}}
\providecommand{\Ad}{}\renewcommand{\Ad}{\operatorname{Ad}}
\providecommand{\Hom}{}\renewcommand{\Hom}{\operatorname{Hom}}
\providecommand{\Iso}{}\renewcommand{\Iso}{\operatorname{Iso}}
\providecommand{\vol}{}\renewcommand{\vol}{\operatorname{vol}}
\providecommand{\supp}{}\renewcommand{\supp}{\operatorname{supp}}
\providecommand{\diag}{}\renewcommand{\diag}{\operatorname{diag}}
\providecommand{\GL}{}\renewcommand{\GL}{\mathrm{GL}}
\providecommand{\SL}{}\renewcommand{\SL}{\mathrm{SL}}
\providecommand{\Gb}{}\renewcommand{\Gb}{\ov{G}}
\providecommand{\Zb}{}\renewcommand{\Zb}{\ov{Z}}
\providecommand{\GDb}{}\renewcommand{\GDb}{\ov{G}_D}
\providecommand{\pib}{}\renewcommand{\pib}{\ov{\pi}}
\providecommand{\ks}{}\renewcommand{\ks}{\kappa^{*}}
\providecommand{\Eb}{}\renewcommand{\Eb}{\ov{E}}
\providecommand{\Tb}{}\renewcommand{\Tb}{\ov{T}}
\providecommand{\cG}{}\renewcommand{\cG}{\mathcal{G}}
\providecommand{\cGb}{}\renewcommand{\cGb}{\ov{\mathcal{G}}}
\providecommand{\iD}{}\renewcommand{\iD}{\iota_D}
\providecommand{\iG}{}\renewcommand{\iG}{\iota_{\GL}}
\providecommand{\jj}{}\renewcommand{\jj}{\jmath}
\providecommand{\Pv}{}\renewcommand{\Pv}{\textbf{[P]}}
\providecommand{\Ev}{}\renewcommand{\Ev}{\textbf{[E]}}
\providecommand{\Ov}{}\renewcommand{\Ov}{\textbf{[O]}}
\providecommand{\Cv}{}\renewcommand{\Cv}{\textbf{[C]}}
\section{Notation specific to \texorpdfstring{Chapter~\ref{ch:I}}{Chapter 4}}\label{I:subsec:glossI}
\begin{description}[style=nextline,leftmargin=1.5em]
\item[$\iD$, $\iG$; $\GDb(K_v)$, $\Gb_v$, $\Zb_v$] The embeddings of $D_v^\times$ and $\GL(r,K_v)$ in $\cG_v=\GL(r,L_v)$, and the inverse images of their images and of $K_v^\times$ in the ambient cover $\cGb_v$ (\S\ref{I:subsec:ambient}, Definition~\ref{I:def:localcovers}).
\item[$\alpha_v$, $\psi_v$] For $v\notin S$: an isomorphism $D_v\cong M_r(K_v)$, and the induced isomorphism of covers $\GDb(K_v)\to\Gb_v$ given by conjugation in $\cGb_v$ (Proposition~\ref{I:prop:outsideS}). $\psi^V=\otimes_{v\notin V}\psi_v$ (\S\ref{I:subsec:sep}).
\item[$E$, $j$, $N$; $\Eb_D$, $\Eb_{\GL}$] At $v\in S$: a maximal subfield $E\subset D_v$, an embedding $j:E\hookrightarrow M_r(F)$, the norm $N=N_{E/F}$, and the covered tori $p^{-1}(\iD(E^\times))$ and $p^{-1}(\iG j(E^\times))$ (\S\ref{I:sec:tori}).
\item[$\Iso_{\Zb}(\Eb_D,\Eb_{\GL})$] Isomorphisms of covered tori over $\iG j\iD^{-1}$ that are the identity on $\mu_n$ and on $\Zb_v$; a torsor under the characters $E^\times\to\mu_n$ trivial on $F^\times$ (Theorem~\ref{I:thm:ident}).
\item[$\mathcal C^{g}(\GDb(K_v))$, $\mathcal C^{g}_{\mathrm{ell}}(\Gb_v)$] The good regular, respectively good regular elliptic, conjugacy classes (\S\ref{I:subsec:cons}).
\item[$E_i$, $j_i$, $\jj_i$; $dg_D$, $dg$, $dt_i$] Representatives of the degree-$r$ extensions of $F$ with fixed embeddings, the corresponding identifications, and the local Haar measures (\S\ref{I:subsec:oi}).
\item[Admissible function; $W_{i,\bullet}$] A locally constant, $\tilde\theta_v^{-1}$-equivariant, $N_{G_\bullet}(T_{i,\bullet})$-invariant function on good regular elements of a covered torus with support compact modulo $\Zb_v$; $W_{i,\bullet}=N_{G_\bullet}(T_{i,\bullet})/T_{i,\bullet}$ (\S\ref{I:subsec:transported}).
\item[$\cGb(\A_L)$, $\GDb(\A_K)$, $\Gb(\A_K)$, $\Zb(\A_K)$] The global ambient cover and its subgroups (\S\ref{I:subsec:globalcovers}).
\item[$\iota$, $\ks_v$, $\kappa_w$, $\tau_w$, $s_N$] The rational splitting, the compact splittings, the Flicker--Kazhdan function $\kappa_w$, the normalized cocycle (Lemma~\ref{I:lem:splittings}), and the unipotent splittings (Proposition~\ref{I:prop:global}(e)).
\item[$S_D$] A finite set of places containing $S$, the archimedean places, the places above $n$, and the places where $L$ or $\tilde\theta$ ramifies or the integrality conditions of Proposition~\ref{I:prop:outsideS}(c) fail (\S\ref{I:subsec:measures}).
\item[$W_\gamma$, $c(\gamma)$, $\mathcal O_H$] For a regular elliptic rational $\gamma$: the automorphisms of $K(\gamma)/K$ fixing $K^\times\gamma$, the class-counting constant, and the set of good regular elliptic rational classes (\S\ref{I:sec:TF}).
\item[$w_0$, $u$; admissible pair] A finite place $w_0\notin S$ carrying the supercusp form, a place $u\in S$ with elliptic support, and the pairs of test functions of Definition~\ref{I:def:globalpairs}.
\item[(A1), (A2); (R1)--(R5)] The analytic inputs of the proof of Proposition~\ref{I:prop:FKvariant}, and the properties of $\Gb(\A_K)$ verified in \S\ref{I:app:cusp}.
\item[$\sigma$, $\tilde f_\sigma$, $d(\sigma)$] A genuine supercuspidal representation of $\Gb_{w_0}$, its normalized matrix coefficient, and its formal degree (Lemma~\ref{I:lem:supercusp}).
\item[$V$, $\rho$] $V=S\cup\{w_0\}$, and $\rho$ an admissible genuine irreducible representation of $\Gb(\A^V)$ (\S\ref{I:subsec:sep}).
\item[Corner Hecke algebra; $e_v$, $e_{\mathfrak t}$, $A_v$, $A_\infty$, $A$, $e\tau$] For an idempotent (or idempotent multiplier) $e$, the algebra $e\mathcal He$; the central-character idempotents, the archimedean $K_\infty$-type projector, the corner algebras and the module $e\tau$ of Lemma~\ref{I:lem:corner}. Irreducible $\tau$ with $e\tau\ne0$ correspond to simple $e\mathcal He$-modules.
\item[Near-equivalence packets $P_D(\rho)$, $P(\rho)$] The automorphic (resp.\ cuspidal) representations with components $\rho$ away from $V$ and $\sigma$ at $w_0$, counted with multiplicity; they need not be finite (Proposition~\ref{I:prop:packet}). $\Omega_D$, $\Omega$ are the finite sets of representations that can contribute at a fixed level (proof of Proposition~\ref{I:prop:packet}).
\item[$\mathcal T_v$, $a$, $\Lambda_a$] The genuine irreducibles of $\GDb(K_v)$ with fixed vectors at a given level, a coefficient function on $\prod_{v\in S}\mathcal T_v$, and the corresponding trace functional (Lemma~\ref{I:lem:cancel}).
\item[$(\mathrm H_D)$, (H); elliptic detectability] The condition that the character at $S$ of the $D$-side (resp.\ $\GL$-side) packet is not identically zero on matching regular-supported test functions (Definition~\ref{I:def:HD}, Theorem~\ref{I:thm:converse}).
\item[(D1)--(D4)] The deferred steps of \S\ref{I:sec:deferred}.
\item[$\Psi$, $\tilde\Psi$, $\Gb_T$, $T^{\reg}$, $\Tb^{g,\reg}$] The conjugation maps, the open set of good elements conjugate into $T^{\reg}$, the regular elements of an elliptic torus $T$, and the good regular elements of $\Tb$ (\S\ref{I:app:oi}).
\end{description}
\endgroup
\begingroup
% macros of Part~II, Draft~10
\providecommand{\ov}{}\renewcommand{\ov}[1]{\overline{#1}}
\providecommand{\A}{}\renewcommand{\A}{\mathbb{A}}
\providecommand{\C}{}\renewcommand{\C}{\mathbb{C}}
\providecommand{\Z}{}\renewcommand{\Z}{\mathbb{Z}}
\providecommand{\cO}{}\renewcommand{\cO}{\mathcal{O}}
\providecommand{\fp}{}\renewcommand{\fp}{\mathfrak{p}}
\providecommand{\gen}{}\renewcommand{\gen}{\mathrm{gen}}
\providecommand{\reg}{}\renewcommand{\reg}{\mathrm{reg}}
\providecommand{\tr}{}\renewcommand{\tr}{\operatorname{tr}}
\providecommand{\Nrd}{}\renewcommand{\Nrd}{\operatorname{Nrd}}
\providecommand{\Ad}{}\renewcommand{\Ad}{\operatorname{Ad}}
\providecommand{\Iso}{}\renewcommand{\Iso}{\operatorname{Iso}}
\providecommand{\vol}{}\renewcommand{\vol}{\operatorname{vol}}
\providecommand{\supp}{}\renewcommand{\supp}{\operatorname{supp}}
\providecommand{\Cent}{}\renewcommand{\Cent}{\operatorname{Cent}}
\providecommand{\GL}{}\renewcommand{\GL}{\mathrm{GL}}
\providecommand{\St}{}\renewcommand{\St}{\mathrm{St}}
\providecommand{\Gb}{}\renewcommand{\Gb}{\ov{G}}
\providecommand{\Zb}{}\renewcommand{\Zb}{\ov{Z}}
\providecommand{\GDb}{}\renewcommand{\GDb}{\ov{G}_D}
\providecommand{\pib}{}\renewcommand{\pib}{\ov{\pi}}
\providecommand{\Eb}{}\renewcommand{\Eb}{\ov{E}}
\providecommand{\Tb}{}\renewcommand{\Tb}{\ov{T}}
\providecommand{\cG}{}\renewcommand{\cG}{\mathcal{G}}
\providecommand{\cGb}{}\renewcommand{\cGb}{\ov{\mathcal{G}}}
\providecommand{\iD}{}\renewcommand{\iD}{\iota_D}
\providecommand{\iG}{}\renewcommand{\iG}{\iota_{\GL}}
\providecommand{\jj}{}\renewcommand{\jj}{\jmath}
\providecommand{\sg}{}\renewcommand{\sg}{\sigma}
\providecommand{\Pv}{}\renewcommand{\Pv}{\textbf{[P]}}
\providecommand{\Ev}{}\renewcommand{\Ev}{\textbf{[E]}}
\providecommand{\Ov}{}\renewcommand{\Ov}{\textbf{[O]}}
\providecommand{\Cv}{}\renewcommand{\Cv}{\textbf{[C]}}
\section{Notation specific to \texorpdfstring{Chapter~\ref{ch:II}}{Chapter 5}}\label{II:subsec:glossII}
\begin{description}[style=nextline,leftmargin=1.5em]
\item[$G$, $\Gb$, $Z$, $\Zb$; $P_g$, $A_g$; eigenvalues] At a finite place $v$ with $F=K_v$: $G$ is $\GL(r,F)$ or $D_v^\times$, $\Gb=p^{-1}(G)$, $Z=F^\times$, $\Zb=p^{-1}(Z)$; $P_g$ is the (reduced) characteristic polynomial of $g$, $A_g=F[g]$, and the eigenvalues of $g$ are the roots of $P_g$ in $\ov F$ (\S\ref{II:subsec:localnotation}).
\item[$\Phi^z$] The orbital integral over $Z\backslash G$, $\Phi^z(\gamma,f)=\vol(T_\gamma/Z)\,\Phi(\gamma,f)$ for regular elliptic $\gamma$ \cite[\S6]{FK} (\S\ref{II:subsec:localnotation}).
\item[$d(\St)$, compatible measures, $\lambda_v$] The formal degree of the Steinberg representation of $\GL(r,F)$; the normalization $\vol(D_v^\times/F^\times)=d(\St)^{-1}$ at $v\in S$ (Definition~\ref{II:def:measures}); and $\lambda_v=d(\St)\vol(D_v^\times/F^\times)$, which equals $1$ for these measures (Remark~\ref{II:rem:measures}).
\item[(I1)--(I11)] The results of Chapter~\ref{ch:I} used here, restated in the form used (\S\ref{II:subsec:fromPartI}).
\item[$\Omega_M$, $V_M$, $V_{D,M}$] The monic degree-$r$ polynomials congruent to $(X-1)^r$ modulo $\fp^M$, and the conjugation-invariant open and closed set of elements of $G$ (resp.\ $D_v^\times$) whose characteristic polynomial lies in $\Omega_M$ (\S\ref{II:subsec:roots}).
\item[$\delta_n$, $B$, $M_0$, $M_1$] The constants of Lemmas~\ref{II:lem:roots} and~\ref{II:lem:chooseM}: $\delta_n=1$ if $p\nmid n$ and $q^{-e/(p-1)}$ if $p\mid n$, $B=|n|_F^{-1}q^{e/(p-1)}$, $M_0$ the level from which the binomial series converges, and $M_1\ge M_0$ the level from which $\tilde\theta\circ\sg$ is trivial on the relevant scalars.
\item[$g^{1/n}$] The canonical $n$-th root $\sum_k\binom{1/n}{k}(g-1)^k\in A_g$ for $g\in V_M$, $M\ge M_0$ (Lemma~\ref{II:lem:roots}).
\item[$\sg$, $\sg_D$, $\sg_{\GL}$] The section $\sg(g)=\tilde h^{\,n}$ with $h=g^{1/n}$ (\S\ref{II:subsec:section}, Proposition~\ref{II:prop:section}).
\item[$\mathcal T$, $\mathcal T_D$; $C_c^\infty(\Gb,\tilde\theta)_{\gen,M}$, $C_c^\infty(G/Z)_M$] The transport $(\mathcal T\tilde f)(p(z)g)=\tilde\theta(z)\tilde f(z\sg(g))$, a bijection between genuine functions supported on $p^{-1}(ZV_M)$ and functions on $G/Z$ supported on $ZV_M$ (Proposition~\ref{II:prop:transport}).
\item[$\Gamma_{\mathcal O}$, $u_{\mathcal O}$] Shalika germs of $\GL(r,F)$ for the fixed measures, indexed by the unipotent classes $\mathcal O$ (partitions of $r$), with representatives $u_{\mathcal O}$; $\Gamma_{\{1\}}$ is the trivial germ (Corollaries~\ref{II:cor:germexp} and~\ref{II:cor:trivgerm}).
\item[(L1), (L2); $f_{\St}$] The two linear inputs: the Steinberg pseudo-coefficient $f_{\St}$ and its orbital integrals, and linear matching of functions with vanishing non-elliptic orbital integrals (\S\ref{II:subsec:linear}).
\item[$e_D$, $\tilde f^{\St}$] The genuine function on $\GDb(K_v)$ equal to $\tilde\theta_v^{-1}$ on $\Zb_v\sg_D(V_{D,M})$, and the transport of $d(\St)^{-1}f_{\St}$ cut off to $ZV_M$; they match (Proposition~\ref{II:prop:eD}).
\item[$\tilde f_{\reg}$] The part of the matching function of Theorem~E supported on the good regular elliptic set (Theorem~\ref{II:thm:matching}).
\item[(Ell), (Non); (G1), (G2)] The orbital-integral form and the unipotent (germ) form of the condition that a function match $\tilde f_D$ near $\Zb_v$ (\S\ref{II:subsec:germ}, Theorem~\ref{II:thm:germ}).
\item[$u$-admissible pair] A pair of global test functions with elliptic support at one place $u\in S$ only and arbitrary matching functions at the other places of $S$ (Definition~\ref{II:def:uadm}).
\item[$(\mathrm H_D')$, $(\mathrm H')$] The weakened non-cancellation hypotheses (Definition~\ref{II:def:HDprime}).
\end{description}
\endgroup
\begingroup
% macros of Part~III, Draft~10
\providecommand{\A}{}\renewcommand{\A}{\mathbb A}
\providecommand{\C}{}\renewcommand{\C}{\mathbb C}
\providecommand{\Z}{}\renewcommand{\Z}{\mathbb Z}
\providecommand{\GL}{}\renewcommand{\GL}{\operatorname{GL}}
\providecommand{\Nrd}{}\renewcommand{\Nrd}{\operatorname{Nrd}}
\providecommand{\Irr}{}\renewcommand{\Irr}{\operatorname{Irr}}
\providecommand{\vol}{}\renewcommand{\vol}{\operatorname{vol}}
\providecommand{\tH}{}\renewcommand{\tH}{\widetilde H}
\providecommand{\tZ}{}\renewcommand{\tZ}{\widetilde Z}
\providecommand{\ttT}{}\renewcommand{\ttT}{\widetilde T}
\providecommand{\reg}{}\renewcommand{\reg}{\mathrm{reg}}
\providecommand{\gen}{}\renewcommand{\gen}{\mathrm{gen}}
\providecommand{\JL}{}\renewcommand{\JL}{\mathrm{JL}}
\providecommand{\Sha}{}\renewcommand{\Sha}{\operatorname{Sha}}
\providecommand{\Gb}{}\renewcommand{\Gb}{\overline{G}}
\providecommand{\GDb}{}\renewcommand{\GDb}{\overline{G}_D}
\providecommand{\Zb}{}\renewcommand{\Zb}{\overline{Z}}
\providecommand{\Ptag}{}\renewcommand{\Ptag}{\textnormal{\textbf{[P]}}}
\providecommand{\Etag}{}\renewcommand{\Etag}{\textnormal{\textbf{[E]}}}
\providecommand{\Otag}{}\renewcommand{\Otag}{\textnormal{\textbf{[O]}}}
\providecommand{\Ctag}{}\renewcommand{\Ctag}{\textnormal{\textbf{[C]}}}
\section{Notation and terminology specific to \texorpdfstring{Chapter~\ref{ch:III}}{Chapter 6}}\label{III:subsec:glossIII}
\begin{description}[style=nextline,leftmargin=2.7cm,labelwidth=2.5cm,font=\normalfont\bfseries]
\item[Adams operation $\psi^n$]
The $\lambda$-ring operation on the representation ring satisfying
$\chi_{\psi^nV}(g)=\chi_V(g^n)$; it can produce a virtual character
rather than an actual representation.
\item[Equivariant $K$-theory]
At a point, $K_Q^0(\mathrm{pt})=R(Q)$, the representation ring
of a finite group $Q$. Exterior powers give its $\lambda$-operations.
\item[Virtual character]
An integral linear combination of irreducible characters.
\item[Hall $\pi$-subgroup]
A subgroup whose order involves only primes in $\pi$ and whose
index is prime to every prime in $\pi$; in the tame proof it
separates the central $n$-part of a finite quotient.
\item[Normalized orbital integral]
Here $J(t,f)=\Delta(t)O(t,f)$, where $\Delta$ is the square root
of the Weyl discriminant and quotient Haar measures are specified.
\item[Transforms $L_D,\mathcal O_D,\mathcal B$]
Respectively the unitary invariant class-function transform, the
forward test-function transfer, and its inverse modulo orbital-integral
kernels; see Definitions~\ref{III:def:L} and Theorem~\ref{III:thm:test-transfer}.
\item[Stable/rational conjugacy]
Conjugacy over an algebraic closure (with descent) versus conjugacy
by the rational points $G(K)$; one stable class may contain several
rational classes for a general reductive group.
\item[Inner form]
A twist of a reductive group by inner automorphisms;
$\GL_1(D)$ is an inner form of $\GL_r$.
\item[Induced torus]
In this paper $T=\operatorname{Res}_{E/K}\mathbb G_m$,
with $T(K)=E^\times$ and $H^1(K,T)=1$ by Hilbert~90.
\item[Kottwitz torus obstruction]
A global cohomological obstruction to transferring an admissible
torus to an inner form despite local embeddings, discussed in
\cite[\S9.5]{Kot84}.
\item[Relative Tamagawa number]
The ratio $\tau_1(G)=\tau(G)/\tau(G_{\mathrm{sc}})$; it is not
a local probability Haar measure.
\item[Brauer class]
The class $[D]\in\operatorname{Br}(K)$ of a central simple algebra;
$E$ embeds maximally in $D$ exactly when $D\otimes_K E$ splits.
\item[Scalar power quotient]
$\mathcal Q(E/K,n)=E^\times/(K^\times E^{\times n})$;
it measures global scalar-modulo-$n$th-power classes.
\item[Kummer and $\Sha^1$]
The exact sequence $1\to T[n]\to T\xrightarrow{n}T\to1$
relates powers to Galois cohomology. $\Sha^1(K,T[n])$ is the
kernel of localization of $H^1(K,T[n])$ at every place.
\item[$H$, $Z$, $\tH$, $\tZ$, $\widetilde T$]
$H=D^\times$ for a local division algebra $D$ of prime degree $r$, $Z=F^\times$, and their inverse images in the cover of Chapter~\ref{ch:I}; $\widetilde T=p^{-1}(T)$ for a maximal torus $T=E^\times$ (\S\ref{III:sec:setting}).
\item[$H^1$, $U_D^1$]
The reduced-norm-one group $\ker(\Nrd)$ and the principal units $1+\mathfrak P_D$ (\S\ref{III:sec:tame}).
\item[Tame, wild]
In this paper \emph{tame} means $p\nmid n$; it imposes no condition on the ramification of maximal subfields (\S\ref{III:sec:tame}).
\item[$\mathcal H_\theta(H)$, $\mathcal H_{\widetilde\theta}(\tH)$]
The Hilbert spaces of $\theta$-equivariant, respectively genuine $\tilde\theta$-equivariant, invariant $L^2$ functions, with the quotient probability measures (\S\ref{III:sec:local}).
\item[$b_D$, $c_D$; character matrix]
$b_D=|n|_F^{-(r-1)/2}$ (Theorem~\ref{III:thm:unitary}) and $c_D=|n|_F^{-(r^2-1)/2}$ (Corollary~\ref{III:cor:integrality}); the matrix of $L_D$ in the bases of irreducible characters (Corollary~\ref{III:cor:matrix}).
\item[$\mathcal C_{\eta^{-1}}(H)$, $L_D^{(\widetilde\eta)}$]
Locally constant functions transforming by $\eta^{-1}$ under $Z$, and the transform of Definition~\ref{III:def:L} for the genuine central character $\widetilde\eta$ (\S\ref{III:sec:orbital}).
\item[$\Psi_D$, (\ref{III:eq:global-conj})]
The conjectural map from genuine irreducible representations to $\Irr(D_v^\times)$ and the target spectral identity of \S\ref{III:sec:global}.
\end{description}
\endgroup
\begingroup
% macros of Part~IV, concluding Draft~6
\providecommand{\F}{}\renewcommand{\F}{F}
\providecommand{\K}{}\renewcommand{\K}{K}
\providecommand{\A}{}\renewcommand{\A}{\mathbb A}
\providecommand{\C}{}\renewcommand{\C}{\mathbb C}
\providecommand{\Z}{}\renewcommand{\Z}{\mathbb Z}
\providecommand{\GL}{}\renewcommand{\GL}{\operatorname{GL}}
\providecommand{\Irr}{}\renewcommand{\Irr}{\operatorname{Irr}}
\providecommand{\Nrd}{}\renewcommand{\Nrd}{\operatorname{Nrd}}
\providecommand{\Nm}{}\renewcommand{\Nm}{\operatorname{N}}
\providecommand{\vol}{}\renewcommand{\vol}{\operatorname{vol}}
\providecommand{\Gal}{}\renewcommand{\Gal}{\operatorname{Gal}}
\providecommand{\Aut}{}\renewcommand{\Aut}{\operatorname{Aut}}
\providecommand{\reg}{}\renewcommand{\reg}{\mathrm{reg}}
\providecommand{\tH}{}\renewcommand{\tH}{\widetilde H}
\providecommand{\tZ}{}\renewcommand{\tZ}{\widetilde Z}
\providecommand{\tT}{}\renewcommand{\tT}{\widetilde T}
\providecommand{\Gb}{}\renewcommand{\Gb}{\overline G}
\providecommand{\GDb}{}\renewcommand{\GDb}{\overline{G_D}}
\providecommand{\Ptag}{}\renewcommand{\Ptag}{\textnormal{\textbf{[P]}}}
\providecommand{\Etag}{}\renewcommand{\Etag}{\textnormal{\textbf{[E]}}}
\providecommand{\Ctag}{}\renewcommand{\Ctag}{\textnormal{\textbf{[C]}}}
\providecommand{\Otag}{}\renewcommand{\Otag}{\textnormal{\textbf{[O]}}}
\section{Terminology specific to \texorpdfstring{Chapter~\ref{ch:IV}}{Chapter 7}}\label{IV:sec:glossary}
\begin{description}[style=nextline,leftmargin=3.0cm]
\item[Good element] A regular lift commuting with the full preimage of its ordinary centralizer; genuine invariant characters vanish on the nongood regular locus.
\item[Scalar-central cover] A cover whose entire inverse image of $F^\times I_r$ is central; here forced by $n\mid r-1+2rc_v$.
\item[Power-map kernel] $\ker(P_n)$ on $T(F)/F^\times$, distinct from the image index.  The kernel can be nontrivial on product tori even when the corresponding field-torus map is injective.
\item[Star map] $h^*=\widetilde h^{\,n}$, independent of lift; multiplication on tori uses $n$ odd.
\item[Weyl discriminant] The factor $\Delta=|\det(1-\operatorname{Ad}(t)|\mathfrak g/\mathfrak t)|^{1/2}$; it is not the unnormalized characteristic-polynomial discriminant.
\item[Toral Fourier test] Expand the discriminant ratio $R_E$ in characters of the compact abelian group $E^\times/F^\times$.  Integral positive coefficients ensure torus-restriction positivity, not full-group integrality.
\item[Nonabelian gluing] The additional requirement that compatible genuine covered-torus characters arise from an integral character of the full group $\widetilde{D^\times}$.
\item[Twisted finite quotient] A finite quotient of $H/Z$ with the central character retained as a multiplier, used to compute projective/genuine irreducible characters.
\item[Packet identity] An identity of traces over collections of automorphic representations; it does not assert a bijection between individual local irreducibles.
\item[Formal degree] The Plancherel normalization of a square-integrable representation relative to a specified quotient Haar measure.
\item[Root-sum relation] Zou's equation (3.2), which sums over all $n$th roots satisfying the stated regularity condition rather than choosing one arbitrarily.
\end{description}
\endgroup
